\documentclass[a4paper,
		12pt,
		parskip=full,
		titlepage
		]{scrartcl}

\usepackage[T1]{fontenc}
\usepackage{graphicx}
\usepackage{geometry}
\usepackage{titlesec}
\usepackage[utf8]{inputenc}
\usepackage{color}
\usepackage{hyperref}
\usepackage{float}


\usepackage{paralist}
\usepackage[section]{placeins} %don't float to next section

\usepackage[lofdepth,lotdepth]{subfig}

%\usepackage{minted}

\hyphenpenalty=10000


%\geometry{a4paper,left=40mm,right=40mm, top=20mm, bottom=40mm}

\linespread{1.25}

\graphicspath{{abbildungen/}} 
\titleformat{\section}[block]{\sffamily\Large\bfseries\filcenter}{\thesection}{1em}{}
\pagestyle{empty}

\newcommand{\blankpage}{
\newpage
\thispagestyle{empty}
\mbox{}
\newpage
}

%opening
\title{Dokument}
\author{}
\date{}

\begin{document}
\thispagestyle{empty}

\hspace{20cm}
\vspace{-3cm}

\begin{figure}[H] \hspace{11cm}
\includegraphics[width=3.2 cm]{HU_Logo}
\end{figure}
% \vspace{5cm}
\begin{center}
  % \vspace{0.5 cm}
  \huge{\bf Title} \\ % Hier fuegen Sie den Titel Ihrer Arbeit ein.
  \vspace{1cm}
  \LARGE  Bachelorarbeit \\ % Geben Sie anstelle der Punkte an, ob es sich um eine
                % Diplomarbeit, eine Masterarbeit oder eine Bachelorarbeit handelt.
  \vspace{1cm}
  \Large zur Erlangung des akademischen Grades \\
  Bachelor of Science (B. Sc.)\\ % Bitte tragen Sie hier anstelle der Punkte ein:
         % Diplominformatiker(in),
         % Bachelor of Arts (B. A.),
         % Bachelor of Science (B. Sc.),
         % Master of Education (M. Ed.) oder
         % Master of Science (M. Sc.).
  \vspace{1.5cm}
  {\large
    \bf{
      \scshape
      Humboldt-Universit\"at zu Berlin \\
      Mathematisch-Naturwissenschaftliche Fakult\"at II \\
      Institut f\"ur Informatik\\
    }
  } 
  % \normalfont
\enlargethispage{10\baselineskip}
\end{center}
\vspace {4 cm}% gegebenenfalls kleiner, falls der Titel der Arbeit sehr lang sein sollte
%{3.2 cm} bei Verwendung von scrreprt, gegebenenfalls kleiner, falls der Titel der Arbeit sehr lang sein sollte
{\large
  \begin{tabular}{llll}
    eingereicht von:    &  && \\ % Bitte Vor- und Nachnamen anstelle der Punkte eintragen.
    geboren am:         &  && \\
    in:                 &  && \\
    &&&\\
    Gutachter: &  && \\
              &  && \\% Bitte Namen der Gutachter(innen) anstelle der Punkte eintragen
                 % bei zwei männlichen Gutachtern kann das (innen) weggestrichen werden
    &&&\\
    eingereicht am:     &   \hspace{3cm} verteidigt am: &  \\ % Bitte lassen Sie
                                    % diese beiden Felder leer.
                                    % Loeschen Sie ggf. das letzte Feld, wenn
                                    % Sie Ihre Arbeit laut Pruefungsordnung nicht
                                    % verteidigen muessen.
  \end{tabular}
}

\newpage

\chapter{Statement of authorship}

\thispagestyle{empty}

%%%%%%%%%%%%%%%%%%%%%%%%%%%%%%%%%%%%%%%%%%%%%%%%%%%%%%%%%%%%%%%%%%%%%%%%%%%%%%%%%%%%%%%%%%%%%%%%%%%%
%% Selbststaendigkeitserklaerung
%%%%%%%%%%%%%%%%%%%%%%%%%%%%%%%%%%%%%%%%%%%%%%%%%%%%%%%%%%%%%%%%%%%%%%%%%%%%%%%%%%%%%%%%%%%%%%%%%%%%

{\parindent 0cm
%%%%%%%%%%%%%%%%%%%%%%%%%%english version%%%%%%%%%%%%%%%%%%%%%%%%%%%%%%
 
I declare that I completed this thesis on my own and that information which has been
directly or indirectly taken from other sources has been noted as such. Neither this
nor a similar work has been presented to an examination committee.

  \vspace{3\baselineskip}
 
  Berlin, \today \hspace{0.25\linewidth}\parbox{0.3\linewidth}{\dotfill}
}
\newpage

\begin{abstract}

\end{abstract}

\addcontentsline{toc}{section}{Contents}
\setcounter{page}{1}
\tableofcontents{}


\pagebreak

\pagestyle{headings}

\newpage
\addcontentsline{toc}{section}{References}
\bibliographystyle{acm}
\bibliography{refs.bib,paper-build/conferences-crossref}

\end{document}
